\documentclass[11pt]{article}
\usepackage{graphicx} % This lets you include figures
\usepackage{hyperref} % This lets you make links to web locations
\usepackage[margin=0.5in]{geometry}
\usepackage[rightcaption]{sidecap}
\usepackage{subcaption}
\usepackage{wrapfig}
\usepackage{float}
\usepackage{imakeidx}
\usepackage{indentfirst}
\makeindex
%---------------------------Do Not Edit Anything Above This Line!!------------------------

% edit the line below, if needed, to change the directory name for your image files.
\graphicspath{ {./images/} }



\begin{document}

%---------------------------Edit Content in the Box to Create the Title Page--------------
\begin{titlepage}
   \begin{center}
       \vspace*{1cm}
	   \Huge
       \textbf{Smart Resume/Job Search App}

       \vspace{0.5cm}
       \Large
       Client: Dr. Galloway \\
       Sprint Number 2 \\
       9/30/2025 \\
   \end{center}

       \vspace{1.5cm}

\begin{table}[h!]
\centering
\begin{tabular}{|l|l|}
\hline
\textbf{Name} & \textbf{Email Address} \\ \hline
Brendan Rucker & brendan.rucker964@topper.wku.edu \\ \hline
Yukai Yang  & yukai.yang457@topper.wku.edu         \\ \hline
Michael Goldsmith & michael.goldsmith370@topper.wku.edu     \\ \hline
\end{tabular}
\end{table}

%Latex Table Generator    
%https://www.tablesgenerator.com/     
        
\vspace{4in}

\centering        
CS 360 \\
Fall 2025\\
Project Organization Documentation

\end{titlepage}
%---------------------------Edit Content in the Box to Create the Title Page--------------


% No text here.


%---------------------------Do Not Edit Anything In This Box!!------------------------
%Table of contents and list of figures will be autogenerated by this section.
\newpage
\setcounter{page}{1}%
\cleardoublepage
\pagenumbering{gobble}
\tableofcontents
\cleardoublepage
\pagenumbering{arabic}
\clearpage
\newpage
\setcounter{page}{1}%
\cleardoublepage
\pagenumbering{gobble}
\listoffigures
\cleardoublepage
\pagenumbering{arabic}
\newpage
%---------------------------Do Not Edit Anything In This Box!!------------------------

% No text here.


%---------------------------Project Team's Organizational Approach------------------------------
\section{Project Team's Organizational Approach} %\section{} is used to create major section headers
%100 words for each sprint
	Sprint 1: %How/where did the group meet?  How often did you meet as an entire team?  Who’s the Project Manager for this sprint?
    Group met in Snell Hall B104 weekly and Discord regularly. We had mandatory meetings with the client weekly and discussed daily through Discord. The project manager for this sprint is Michael.
	
	Sprint 2: %How/where did the group meet?  How often did you meet as an entire team?  Who’s the Project Manager for this sprint?
    We met as a group at least every week, having client meetings every Thursday in Snell Hall B104 @12:35pm. We discuss the project before class on Tuesday and Thursday as well as engage in project related discussion in class when prompted. We also have weekly meetings online on Discord every Wednesday to discuss preparations for our Client Meeting the next day. The project manager for this sprint was Brendan Rucker. 

	Sprint 3: %How/where did the group meet?  How often did you meet as an entire team?  Who’s the Project Manager for this sprint?

	Sprint 4: %How/where did the group meet?  How often did you meet as an entire team?  Who’s the Project Manager for this sprint?



%---------------------------End Project Team's Organizational Approach------------------------------


% No text here.


%---------------------------Schedule Organization---------------------------------------------------
\section{Schedule Organization}
%Gantt charts cover the tasks/time commitments and estimations for the entire project.  We will have four iterations of the Gantt Chart, with iteration focusing on a specific sprint.

\subsection{Gantt Chart v1:}
%100 words to describe the focus for this sprint.
%Identify the location for the Gantt Chart created during Sprint 1.  Should be clearly labeled in the project directory.
In GitHub under documentations, focus of this sprint is documentation and research. Creating an outline for future work.



\subsection{Gantt Chart v2:}
%100 words to describe the focus for this sprint.
%Identify the location for the Gantt Chart created during Sprint 2.  Should be clearly labeled in the project directory.
Our Gantt Chart is available in the root of our project directory as well as the root of our file directory on Overleaf. This chart covers who is assigned to which UML Diagram and generally a deadline to get complete the diagram. These deadlines are decided based on when we've discussed the diagrams in class.



\subsection{Gantt Chart v3:}
%100 words to describe the focus for this sprint.
%Identify the location for the Gantt Chart created during Sprint 3.  Should be clearly labeled in the project directory.
Text goes here.



\subsection{Final Gantt Chart:}
%100 words to describe the focus for this sprint.
%Identify the location for the Gantt Chart created during Sprint 4.  Should be clearly labeled in the project directory.
Text goes here.



%---------------------------End Schedule Organization---------------------------------------------------


% No text here.


%---------------------------Progress Visibility---------------------------------------------------
\section{Progress Visibility}
%100 words for each sprint.
%For each sprint, explain how each member of the group is progressing with assigned tasks and how that progress is shared with the group.  Also, explain how the group is progressing with assigned tasks and how that progress is shared with the client.  Examples:  how does the group assign tasks?  How to group members know tasks assigned to them?  How do group members communicate when assigned tasks are complete, need assistance, or waiting on other tasks to be completed first?
\subsection{Sprint 1 Progress Visibility}
We are progressing smoothly, this sprint we mainly focused on research and creating our outline for what we want the application to do and what it will look like. We are looking at language models and deciding which would work the best. Michael suggested that we use two, one model for initial analysis and another for revision and feedback.

\subsection{Sprint 2 Progress Visibility}
There was a bit more struggle with understanding how to build UML diagrams and how to apply those diagrams to our project idea, but at the end I think it we have a broader understanding of the need for these diagrams and are ready to start working on Sprint 3 of applying these diagrams. This sprint relied on heavy communication with the team to ensure everyone understands the proposed idea on how the project will run so that we could work together to create UML Diagrams that build upon the same components instead of adding conflicting components. For this sprint, we assigned tasks as equal as possible based on each members strengths with Brendan working on the Use Case, State, and Structural Class Diagrams, Michael working on the Sequence and Creational Class Diagrams, and Yukai working on the Component, Deployment, and Behavioral Class Diagrams. 

\subsection{Sprint 3 Progress Visibility}
Text goes here.

\subsection{Sprint 4 Progress Visibility}
Text goes here.

%---------------------------End Progress Visibility---------------------------------------------------

% No text here.


%---------------------------Software Process Model---------------------------------------------------
\section{Software Process Model}
%150 words
%Describe in this section the Software Process Model used and how it increases the quality of the final deliverables.  The team should also define the quality control steps that are used in the Software Process Model.
We will be using a modified Waterfall Software Process Model that will split this project into 4 sprints, each sprint focused on a different process. Sprint 1 will focus on the feasibility study, denoting functional and non-functional requirements to come up with a development plan. Sprint 2 is focused on expanding on this idea by creating UML Diagrams that structure what our code could look like. Sprint 3 will focus on the development of the project; building the code. Sprint 4 will focus on testing the project, ensuring all the functions are running correctly and finding any areas of improvement. 

%---------------------------End Software Process Model---------------------------------------------------

% No text here.


%---------------------------Risk Management--------------------------------------------------------------
\section{Risk Management}
%Use this section to describe the team's approach to risk management.

\subsection{Risk Identification}
%List, categorize, and prioritize all potential risks associated with the project.
\textbf{High Priority}

-AI model performance may not meet the 5-second processing requirement.

-Integration with third-party job board APIs may break due to version changes.

-Security vulnerabilities could expose sensitive resume data.

-Scope creep from adding extra features beyond requirements.

-Insufficient access to premium APIs or job board data.
\textbf{Medium Priority}

-Cross-platform compatibility issues (Windows, macOS, Linux).

-Database performance bottlenecks when handling large datasets.

-Limited computing resources for training or testing AI models.
\textbf{Low Priority}

-Limited offline functionality if local datasets are incomplete.

\subsection{Risk Planning}
%Give overviews of plans for specific risk types
Text goes here.

\subsection{Risk Monitoring}
%Give overviews of how the team will monitor specific types of risks
The team will address technical risks by implementing robust testing and monitoring. AI models will be benchmarked regularly to ensure sub-five-second processing.

The team will identify early which APIs are free and which require subscriptions, budgeting accordingly. Development will use lightweight datasets for testing before scaling to larger sets. Team members will ensure they have machines that meet hardware requirements, and cloud resources may be considered if local capacity is insufficient. Collaboration tools will have backup options, such as OneDrive for documents if Google Drive fails.

%---------------------------Risk Management--------------------------------------------------------------





%example image:  uncomment to show usage
%\begin{figure}[h]
%    \centering
%    \includegraphics[width=1\textwidth]{images/Add_non-music.png}
%    \caption{This is how you add non-music items.}
%    \label{fig16}
%\end{figure}


%example links:  uncomment to show usage.
%\url{https://www.youtube.com}
%\href{https://www.wku.edu/}{WKU Homepage}
%\footnote{You can put the link in a footnote like this.}

% Anything to the right of a percent sign will be ignored by LaTeX.
% You can use this to put notes to yourself.  



\end{document}
